% Level 3, Week 4 -- Searching and How Fast Is Fast? (student workbook)
\chapterweek{4}{Searching and How Fast Is Fast?}{Smart algorithms do less work: halving beats checking one by one.}

\section{What you will learn}
\begin{checklist}
  \item Search a list in two ways: linear search and binary search.
  \item Count the steps an algorithm takes, and compare how they grow.
  \item Build a bot that always guesses your number in 7 tries or fewer.
\end{checklist}

\section{Two ways to search}
\begin{analogy}[The guessing game]
Your friend thinks of a number from 1 to 100 and says \enquote{higher} or \enquote{lower} after each guess. If you guess 1, 2, 3, 4\ldots{} you might need 100 guesses. If you guess 50 first, then the middle of what is left, you never need more than 7. That second way is binary search.
\end{analogy}

\textbf{Linear search} looks at the items one by one, from the start. It works on any list, but a long list means a long search.

\textbf{Binary search} only works on a \textbf{sorted} list. It looks at the middle item. If that is too small, the answer must be on the right, so it throws away the whole left half (and the other way round). Every step cuts the list in half.

\begin{center}
%% alt: Binary search in a sorted list of eight numbers: 3, 8, 15, 21, 29, 37, 44 and 52 at positions 0 to 7. Step 1: the middle is 21, which is too small, so throw away the left half. Step 2: the middle is 37, which is too small, so throw away the left half. Step 3: the middle is 44. Found 44!
\begin{tikzpicture}[cell/.style={draw=black,minimum width=0.95cm,minimum height=0.7cm,font=\small\ttfamily},x=1cm,y=1cm]
  \def\vals{{3,8,15,21,29,37,44,52}}
  \foreach \row/\lo/\hi/\mid/\note in {0/0/7/3/{step 1: 21 is too small, throw away the left half},
                                        1/4/7/5/{step 2: 37 is too small, throw away the left half},
                                        2/6/7/6/{step 3: found 44!}}{
    \foreach \i in {0,...,7}{
      \pgfmathtruncatemacro{\v}{\vals[\i]}
      \ifnum\i<\lo
        \node[cell,fill=kcgray!25,text=kcgray] at (\i,-1.3*\row) {\v};
      \else\ifnum\i>\hi
        \node[cell,fill=kcgray!25,text=kcgray] at (\i,-1.3*\row) {\v};
      \else\ifnum\i=\mid
        \node[cell,fill=kcorange!45,line width=1.4pt] at (\i,-1.3*\row) {\v};
      \else
        \node[cell,fill=kcsky!20] at (\i,-1.3*\row) {\v};
      \fi\fi\fi
    }
    \node[font=\scriptsize,above] at (\mid,-1.3*\row+0.3) {mid};
    \node[anchor=west,font=\small,text width=5.2cm] at (8.0,-1.3*\row) {\note};
  }
  \foreach \i in {0,...,7} \node[font=\scriptsize,text=kcgray] at (\i,-3.15) {\i};
  \node[font=\scriptsize,text=kcgray,anchor=east] at (-0.6,-3.15) {index};
\end{tikzpicture}
\end{center}
Gray boxes have been thrown away; the thick orange box is the one we look at (\texttt{mid}). Searching for 44 took 3 steps. Linear search would take 7.

\section{How fast is fast?}
To compare algorithms, count the \textbf{steps} in the \textbf{worst case} (the unluckiest search). For linear search that is quick to count: with 10 items, the unluckiest search looks at all 10.

\begin{tip}[Math helper: how many times can you cut it in half?]
Every binary search step cuts the list in half. Leftover halves are rounded \emph{down} (half of 5 is 2, with 1 left over). Keep halving until nothing is left, and count the numbers you wrote. That count is the worst case for binary search.
\begin{center}
\renewcommand{\arraystretch}{1.25}
\begin{tabular}{@{}rlr@{}}
\toprule
\textbf{Items} & \textbf{Halving ladder} & \textbf{Steps} \\ \midrule
10 & 10, 5, 2, 1 & 4 \\
100 & 100, 50, 25, 12, 6, 3, 1 & 7 \\
1000 & 1000, 500, 250, 125, 62, 31, 15, 7, 3, 1 & 10 \\
\bottomrule
\end{tabular}
\end{center}
Another way to see it: every extra step lets binary search handle a list about \textbf{twice} as big.
\begin{center}
\renewcommand{\arraystretch}{1.2}
\begin{tabular}{@{}lcccccccccc@{}}
\toprule
\textbf{Steps} & 1 & 2 & 3 & 4 & 5 & 6 & 7 & 8 & 9 & 10 \\
\textbf{Biggest list} & 1 & 3 & 7 & 15 & 31 & 63 & 127 & 255 & 511 & 1023 \\
\bottomrule
\end{tabular}
\end{center}
So 100 items need 7 steps (100 fits under 127), and 1000 items need 10 steps (1000 fits under 1023).
\end{tip}

\begin{center}
\begin{tabular}{@{}rrr@{}}
\toprule
\textbf{Items in the list} & \textbf{Linear search} & \textbf{Binary search} \\ \midrule
10 & 10 steps & 4 steps \\
100 & 100 steps & 7 steps \\
1000 & 1000 steps & 10 steps \\
\bottomrule
\end{tabular}
\end{center}
When the list gets 10 times bigger, linear search does 10 times more work. Binary search only does about 3 more steps!

\section{Binary search in Python}
\begin{pycode}
low = 0
high = len(items) - 1          # the LAST index, not the length
while low <= high:             # <= : one item left is still worth checking
    mid = (low + high) // 2    # // makes a whole number
    if items[mid] == target:
        return mid
    elif items[mid] < target:
        low = mid + 1          # throw away the left half (and mid)
    else:
        high = mid - 1         # throw away the right half (and mid)
return -1                      # nothing left: not in the list
\end{pycode}
\textbf{New symbol:} \texttt{//} divides and keeps only the whole part. \texttt{7 // 2} is \texttt{3}, because 7 divided by 2 is 3 remainder 1, and \texttt{//} throws the remainder away. A list place must be a whole number. (\texttt{7 / 2} would give \texttt{3.5}.)\\
\textbf{What is \texttt{-1}?} It is \enquote{minus one}, a number below zero. List places start at 0, so the search uses \texttt{-1} as a signal that means \enquote{not found}.\\
\textbf{\texttt{<=}} means \enquote{less than or equal to}. When \texttt{low} and \texttt{high} are equal, there is still one item left to check.

\begin{tip}[Your turn: trace it before you code it]
Search for \textbf{52} in the list from the picture: \texttt{[3, 8, 15, 21, 29, 37, 44, 52]}. Fill in one row per step until you find it. Step 1 is done for you.
\begin{center}
\renewcommand{\arraystretch}{1.35}
\begin{tabular}{|c|c|c|c|c|p{4.2cm}|}
\hline
\textbf{step} & \textbf{low} & \textbf{high} & \textbf{mid} & \textbf{items[mid]} & \textbf{so\ldots} \\ \hline
1 & 0 & 7 & 3 & 21 & too small: low = 4 \\ \hline
2 & & & & & \\ \hline
3 & & & & & \\ \hline
4 & & & & & \\ \hline
\end{tabular}
\end{center}
How many steps did binary search take? \rule{1.5cm}{0.4pt}\quad Linear search? \rule{1.5cm}{0.4pt}
\end{tip}

\begin{exercise}[Search Race]
Folder: \texttt{week04/ex-search-race}. \texttt{linear\_search} is done. Write \texttt{binary\_search}: it returns \texttt{(index, steps)}. Run it to see the race table.\\
Two values in round parentheses are a \textbf{tuple}. \texttt{return mid, steps} gives back both at once, and \texttt{index, steps = binary\_search(items, 44)} catches both.\\
\textbf{Predict} the binary steps for 10000 items (make a halving ladder: 10000, 5000, 2500, \ldots): \rule{2cm}{0.4pt}\quad What did the program say? \rule{2cm}{0.4pt}\\
\textbf{Done looks like:} the race table ends with \texttt{10000 | 10000 | ...} and \texttt{check\_my\_work.py} passes.
\end{exercise}

\begin{exercise}[Guess-the-Number Bot]
Folder: \texttt{week04/ex-guess-bot}. First run the slow bot: it guesses 1, 2, 3\ldots{} (press the red Stop button). Then:
\begin{itemize}[itemsep=1pt]
  \item finish \texttt{check(guess, secret)}: return \texttt{"higher"}, \texttt{"lower"} or \texttt{"correct"};
  \item make \texttt{next\_guess(low, high)} return the middle.
\end{itemize}
Play 3 games. How many guesses did the bot need? \rule{1.5cm}{0.4pt}\ \rule{1.5cm}{0.4pt}\ \rule{1.5cm}{0.4pt}\\
\textbf{Done looks like:} the first line says \texttt{Bot practice: secret 73 found with guesses [50, 75, 62, 68, 71, 73]}, and the bot never needs more than 7 guesses.
\end{exercise}

\begin{tip}[Filter bubbles]
Search engines are amazingly fast, but they also choose which results to show you first, based on what you clicked before. If you only ever see what you already like, you are in a \enquote{filter bubble}. Try different search words, look at more than one source, and ask a grown-up when something seems odd.
\end{tip}

\begin{vocab}
  \vocabword{algorithm}{A clear list of steps that solves a problem.}
  \vocabword{linear search}{Check items one by one from the start.}
  \vocabword{binary search}{In a sorted list, check the middle and throw away the half that cannot hold the answer.}
  \vocabword{worst case}{The unluckiest situation, when an algorithm needs the most steps.}
  \vocabword{off-by-one error}{A bug where a number is 1 too big or 1 too small, like using \texttt{<} instead of \texttt{<=}.}
\end{vocab}

\begin{reflection}
Explain to a younger friend why binary search needs a sorted list:
\answerlines{2}
Where in real life do you use halving to find something?
\answerlines{1}
\end{reflection}
